\section{Discussion and Limitations}
The proposed benchmark treats \vzd{} as a setting for studying components of general agency: integrating observations, acting toward goals, using memory, and composing behavior over time. Completion of a finite collection of Doom tasks would provide evidence about those capabilities in this domain. It would not establish general intelligence or broad transfer outside the game. The empirical evaluation must determine whether the tasks discriminate among current systems and which differences are stable across runs.

Public Doom levels may be familiar from model pretraining, walkthroughs, or previous benchmarks. Their value as interpretable reference tasks should therefore be separated from claims about generalization. Even with unfamiliar layouts, an agent may benefit from prior knowledge of game mechanics. Future evaluation should distinguish this useful prior knowledge from memorized routes, while using independently authored or generated maps to test adaptation.

Observation access also shapes the task. Automaps, game variables, text notifications, and perfect semantic categories can remove uncertainty that would otherwise require perception or memory. Such assistance is a useful experimental intervention when declared and controlled. Comparisons across different observation privileges, action timing, or training regimes cannot be reduced to a single ordering of agent capability.

The breadth of the library exceeds the scope of the initial MCP benchmark. Audio and notification tools still need to be integrated into that interface, and the PettingZoo wrapper remains a separate development effort. Their evaluation should follow the validation of each observation and task contract. More generally, synchronous gameplay favors agents that can deliberate between actions; a later real-time study would test an additional capability under latency constraints.

Finally, a sparse outcome alone rarely explains failure. The combination of task-level measures and interaction traces is intended to make negative results informative, while keeping diagnostic information separate from what the agent receives. The current manuscript provides that evaluation structure. Claims about performance, modality benefits, or the value of a particular harness remain contingent on the planned experiments.

\section{Conclusion}
\vzd{} has evolved from a primarily visual control platform toward a testbed for multimodal, extended, and social interaction. Audio, textual feedback, semantic categories, full-level tasks, and PettingZoo integration broaden its research scope. The MCP server adds a common route for language-model agents to observe, act, and maintain selective visual memories. Together, these developments support an evaluation that connects standard scenario skills with the opening levels of Doom and Doom~II, while distinguishing the effects of models, harnesses, and observation access. The proposed study and reserved results tables provide the basis for assessing that contribution empirically.
